\chapter{Numerical Methods}\label{chapter:4}

In this chapter we describe and validate the numerical techniques used to solve the optimization problem. In addition, we determine the parameters including the resolution, time step size, regularization parameter $\alpha$, and Gevrey parameters $s$ and $\sigma$.

\section{Discretization}
The Euler-Voigt (\ref{eq:Euler:Voigt}) and adjoint system (\ref{eq:euler:voigt:adjoint}) were discretized in space using a standard pseudospectral Fourier method on a uniform grid with $N=128^3$, $256^3$, $512^3$, $1024^3$ grid points. Derivatives are evaluated in the Fourier space and nonlinear products are computed in the physical space. The standard 2/3 dealiasing rule is also applied.
\\~\\
Time stepping is performed with the explicit fourth-order Runge-Kutta (RK4) method. The determination of our time step size $\Delta t$ is based on two factors. First, we consider the accuracy of the solution, determined by the relative error of the objective functional (\ref{eq:objective:functional}) relative to the results obtained with a minimum step size for a time interval $T=5$. Second, we consider the conservation of the $\alpha$-energy (\ref{eq:alpha:Energy}), determined by the maximum relative error of $E_\alpha(t)$ relative to $E_\alpha(0)$, for $t\in[0,5]$. These factors are analysed for each resolution and value of $\alpha$ for the TGV.
\\~\\
The goal of our optimization search is to identify initial conditions $\beat$ that produce 'extreme' behaviour in solutions to the Euler-Voigt equations. Therefore, the TGV as the initial guess will by definition lead to a less extreme flow evolution than the initial conditions constructed at each iteration of the optimization search. To ensure stability, accuracy, and conservation of $\alpha$-energy upon optimization, we must select a step size smaller than the minimum step size identified by the above criteria.

\subsection{Numerical Stability \& Accuracy}

We begin by analysing the accuracy of the time discretization. To determine the accuracy of the discretization with a given time step size we solve the Euler-Voigt system for $T=5$ with the TGV used as the initial condition and evaluate the objective functional (\ref{eq:objective:functional}). First, we do this for $\Delta t = 2^{-8}$ and designate this result as our 'correct' value $\Phi_c$. We then compute the relative error of the objective functional for each value of $\Delta t$, denoted $\Phi_{\Delta t}$, to $\Phi_{\min}$
\begin{equation}\label{eq:epsilon:phi}
\epsilon_{\Phi} = \left|\frac{\Phi_{\min} - \Phi_{\Delta t}}{\Phi_{\min}}\right|.
\end{equation}

\noindent In \cite{larios_computational_2018}, the regularization parameter was in the range $\alpha=12/1024$, ..., $36/1024$. We extend this range in both directions and use the regularization parameters $\alpha = 8/1024$, ..., $64/1024$. Table \ref{tbl:accuracy} shows the values for $\epsilon_{\Phi}$ for $\alpha=16/1024,...,64/1024$. We include only these values of $\alpha$ in this table because they led to identical results for all four resolutions. This is due to the fact that the regularization term in (\ref{eq:Euler:Voigt}) kept the solution well-resolved even at the lowest resolution, and so the accuracy did not increase with the higher resolutions. Table \ref{tbl:accuracy:2} shows the results of $\epsilon_{\Phi}$ for $\alpha=8/1024,12/1024$. For these values the model becomes under-resolved at the lowest resolution and so accuracy does increase at the finer resolutions. 

\renewcommand{\arraystretch}{1.5}
\begin{table}[]
\centering
\begin{tabular}{lllllll}
                                &                               &                                                       &                                                       & $\alpha$                                              &                                                       &                                                       \\ \cline{3-7} 
                                & \multicolumn{1}{l|}{}         & \multicolumn{1}{l|}{16/1024}                          & \multicolumn{1}{l|}{24/1024}                          & \multicolumn{1}{l|}{32/1024}                          & \multicolumn{1}{l|}{48/1024}                          & \multicolumn{1}{l|}{64/1024}                          \\ \cline{2-7} 
\multicolumn{1}{l|}{}           & \multicolumn{1}{l|}{$2^{-2}$} & \multicolumn{1}{l|}{\cellcolor[HTML]{C0C0C0}7.00E-05} & \multicolumn{1}{l|}{\cellcolor[HTML]{C0C0C0}1.45E-05} & \multicolumn{1}{l|}{\cellcolor[HTML]{C0C0C0}4.80E-06} & \multicolumn{1}{l|}{\cellcolor[HTML]{C0C0C0}8.70E-07} & \multicolumn{1}{l|}{\cellcolor[HTML]{C0C0C0}1.58E-07} \\ \cline{2-7} 
\multicolumn{1}{l|}{}           & \multicolumn{1}{l|}{$2^{-3}$} & \multicolumn{1}{l|}{3.64E-06}                         & \multicolumn{1}{l|}{7.89E-07}                         & \multicolumn{1}{l|}{2.63E-07}                         & \multicolumn{1}{l|}{5.03E-08}                         & \multicolumn{1}{l|}{9.29E-09}                         \\ \cline{2-7} 
\multicolumn{1}{l|}{}           & \multicolumn{1}{l|}{$2^{-4}$} & \multicolumn{1}{l|}{\cellcolor[HTML]{C0C0C0}2.02E-07} & \multicolumn{1}{l|}{\cellcolor[HTML]{C0C0C0}4.53E-08} & \multicolumn{1}{l|}{\cellcolor[HTML]{C0C0C0}1.52E-08} & \multicolumn{1}{l|}{\cellcolor[HTML]{C0C0C0}3.01E-09} & \multicolumn{1}{l|}{\cellcolor[HTML]{C0C0C0}5.63E-10} \\ \cline{2-7} 
\multicolumn{1}{l|}{$\Delta t$} & \multicolumn{1}{l|}{$2^{-5}$} & \multicolumn{1}{l|}{1.21E-08}                         & \multicolumn{1}{l|}{2.70E-09}                         & \multicolumn{1}{l|}{9.16E-10}                         & \multicolumn{1}{l|}{1.83E-10}                         & \multicolumn{1}{l|}{3.48E-11}                         \\ \cline{2-7} 
\multicolumn{1}{l|}{}           & \multicolumn{1}{l|}{$2^{-6}$} & \multicolumn{1}{l|}{\cellcolor[HTML]{C0C0C0}8.63E-10} & \multicolumn{1}{l|}{\cellcolor[HTML]{C0C0C0}1.65E-10} & \multicolumn{1}{l|}{\cellcolor[HTML]{C0C0C0}5.99E-11} & \multicolumn{1}{l|}{\cellcolor[HTML]{C0C0C0}8.33E-12} & \multicolumn{1}{l|}{\cellcolor[HTML]{C0C0C0}2.05E-12} \\ \cline{2-7} 
\multicolumn{1}{l|}{}           & \multicolumn{1}{l|}{$2^{-7}$} & \multicolumn{1}{l|}{9.44E-11}                         & \multicolumn{1}{l|}{9.42E-12}                         & \multicolumn{1}{l|}{5.99E-12}                         & \multicolumn{1}{l|}{0}                                & \multicolumn{1}{l|}{0}                                \\ \cline{2-7} 
%\multicolumn{1}{l|}{}           & \multicolumn{1}{l|}{$2^{-8}$} & \multicolumn{1}{l|}{\cellcolor[HTML]{C0C0C0}0}        & \multicolumn{1}{l|}{\cellcolor[HTML]{C0C0C0}0}        & \multicolumn{1}{l|}{\cellcolor[HTML]{C0C0C0}0}        & \multicolumn{1}{l|}{\cellcolor[HTML]{C0C0C0}0}        & \multicolumn{1}{l|}{\cellcolor[HTML]{C0C0C0}0}        \\ \cline{2-7} 
\end{tabular}
\caption{Results $\epsilon_{\Phi}$ for $\alpha=16/1024,...,64/1024$.}
\label{tbl:accuracy}
\end{table}


\begin{table}[!htb]
	\centering
    \begin{subtable}{.5\linewidth}
      \centering
        \begin{tabular}{llll}
                                &                               & $\alpha$                                              &                                                       \\ \cline{3-4} 
                                & \multicolumn{1}{l|}{}         & \multicolumn{1}{l|}{8/1024}                           & \multicolumn{1}{l|}{12/1024}                          \\ \cline{2-4} 
\multicolumn{1}{l|}{}           & \multicolumn{1}{l|}{$2^{-2}$} & \multicolumn{1}{l|}{\cellcolor[HTML]{C0C0C0}0.0045}   & \multicolumn{1}{l|}{\cellcolor[HTML]{C0C0C0}2.66E-04} \\ \cline{2-4} 
\multicolumn{1}{l|}{}           & \multicolumn{1}{l|}{$2^{-3}$} & \multicolumn{1}{l|}{0.0011}                           & \multicolumn{1}{l|}{1.46E-05}                         \\ \cline{2-4} 
\multicolumn{1}{l|}{}           & \multicolumn{1}{l|}{$2^{-4}$} & \multicolumn{1}{l|}{\cellcolor[HTML]{C0C0C0}4.85E-04} & \multicolumn{1}{l|}{\cellcolor[HTML]{C0C0C0}1.30E-06} \\ \cline{2-4} 
\multicolumn{1}{l|}{$\Delta t$} & \multicolumn{1}{l|}{$2^{-5}$} & \multicolumn{1}{l|}{2.33E-04}                         & \multicolumn{1}{l|}{3.19E-07}                         \\ \cline{2-4} 
\multicolumn{1}{l|}{}           & \multicolumn{1}{l|}{$2^{-6}$} & \multicolumn{1}{l|}{\cellcolor[HTML]{C0C0C0}1.02E-04} & \multicolumn{1}{l|}{\cellcolor[HTML]{C0C0C0}1.23E-07} \\ \cline{2-4} 
\multicolumn{1}{l|}{}           & \multicolumn{1}{l|}{$2^{-7}$} & \multicolumn{1}{l|}{3.45E-05}                         & \multicolumn{1}{l|}{4.05E-08}                         \\ \cline{2-4} 
%\multicolumn{1}{l|}{}           & \multicolumn{1}{l|}{$2^{-8}$} & \multicolumn{1}{l|}{\cellcolor[HTML]{C0C0C0}0}        & \multicolumn{1}{l|}{\cellcolor[HTML]{C0C0C0}0}        \\ \cline{2-4} 
\end{tabular}

    \caption{$N=128^3$}
    \end{subtable}%
    \begin{subtable}{.5\linewidth}
      \centering
      \begin{tabular}{llll}
                                &                               & $\alpha$                                              &                                                       \\ \cline{3-4} 
                                & \multicolumn{1}{l|}{}         & \multicolumn{1}{l|}{8/1024}                           & \multicolumn{1}{l|}{12/1024}                          \\ \cline{2-4} 
\multicolumn{1}{l|}{}           & \multicolumn{1}{l|}{$2^{-2}$} & \multicolumn{1}{l|}{\cellcolor[HTML]{C0C0C0}0.003}    & \multicolumn{1}{l|}{\cellcolor[HTML]{C0C0C0}2.64E-04} \\ \cline{2-4} 
\multicolumn{1}{l|}{}           & \multicolumn{1}{l|}{$2^{-3}$} & \multicolumn{1}{l|}{1.82E-04}                         & \multicolumn{1}{l|}{1.34E-05}                         \\ \cline{2-4} 
\multicolumn{1}{l|}{}           & \multicolumn{1}{l|}{$2^{-4}$} & \multicolumn{1}{l|}{\cellcolor[HTML]{C0C0C0}9.46E-06} & \multicolumn{1}{l|}{\cellcolor[HTML]{C0C0C0}7.05E-07} \\ \cline{2-4} 
\multicolumn{1}{l|}{$\Delta t$} & \multicolumn{1}{l|}{$2^{-5}$} & \multicolumn{1}{l|}{5.12E-07}                         & \multicolumn{1}{l|}{3.95E-08}                         \\ \cline{2-4} 
\multicolumn{1}{l|}{}           & \multicolumn{1}{l|}{$2^{-6}$} & \multicolumn{1}{l|}{\cellcolor[HTML]{C0C0C0}3.00E-08} & \multicolumn{1}{l|}{\cellcolor[HTML]{C0C0C0}2.32E-09} \\ \cline{2-4} 
\multicolumn{1}{l|}{}           & \multicolumn{1}{l|}{$2^{-7}$} & \multicolumn{1}{l|}{1.96E-09}                         & \multicolumn{1}{l|}{1.31E-10}                         \\ \cline{2-4} 
%\multicolumn{1}{l|}{}           & \multicolumn{1}{l|}{$2^{-8}$} & \multicolumn{1}{l|}{\cellcolor[HTML]{C0C0C0}0}        & \multicolumn{1}{l|}{\cellcolor[HTML]{C0C0C0}0}        \\ \cline{2-4} 
\end{tabular}
	
    \caption{$N=256^3,512^3,1024^3$}
    \end{subtable} 
    
    \caption{Results $\epsilon_{\Phi}$ for $\alpha=8/1024,12/1024$.}
    \label{tbl:accuracy:2}
\end{table}

\noindent We now examine the conservation of $\alpha$-energy for the Euler-Voigt model. To determine the accuracy of discretization with a given time step size we solve the Euler-Voigt system for $T=5$ with the TGV used as the initial condition and evaluate the maximum relative error of the $\alpha$-energy
\begin{equation}\label{eq:alpha:energy:rel}
\epsilon_{\alpha} = \max_{t\in[0,T]}\left|\frac{E_\alpha(t) - E_\alpha(0)}{E_\alpha(0)}\right|.
\end{equation}
Table \ref{tbl:alpha:energy} shows the values of $\epsilon_{\alpha}$ for $\alpha=16/1024,...,64/1024$. Just as with the values of $\epsilon_{\Phi}$, we include only these values of $\alpha$ because they led to identical values for all four resolutions. Table \ref{tbl:alpha:energy:2} shows the values of $\epsilon_{\alpha}$ for $\alpha=8/1024,12/1024$. The time step size of each model run was selected based on the results in these tables to ensure accuracy and conservation of the $\alpha$-energy for the chosen value of $\alpha$ and resolution.

\begin{table}[]
\centering
\begin{tabular}{lllllll}
                                &                               &                                                       &                                                       & $\alpha$                                              &                                                       &                                                       \\ \cline{3-7} 
                                & \multicolumn{1}{l|}{}         & \multicolumn{1}{l|}{16/1024}                          & \multicolumn{1}{l|}{24/1024}                          & \multicolumn{1}{l|}{32/1024}                          & \multicolumn{1}{l|}{48/1024}                          & \multicolumn{1}{l|}{64/1024}                          \\ \cline{2-7} 
\multicolumn{1}{l|}{}           & \multicolumn{1}{l|}{$2^{-2}$} & \multicolumn{1}{l|}{\cellcolor[HTML]{C0C0C0}8.73E-06} & \multicolumn{1}{l|}{\cellcolor[HTML]{C0C0C0}1.90E-06} & \multicolumn{1}{l|}{\cellcolor[HTML]{C0C0C0}5.95E-07} & \multicolumn{1}{l|}{\cellcolor[HTML]{C0C0C0}9.75E-08} & \multicolumn{1}{l|}{\cellcolor[HTML]{C0C0C0}1.64E-08} \\ \cline{2-7} 
\multicolumn{1}{l|}{}           & \multicolumn{1}{l|}{$2^{-3}$} & \multicolumn{1}{l|}{\cellcolor[HTML]{FFFFFF}3.72E-07} & \multicolumn{1}{l|}{\cellcolor[HTML]{FFFFFF}7.68E-08} & \multicolumn{1}{l|}{\cellcolor[HTML]{FFFFFF}2.29E-08} & \multicolumn{1}{l|}{\cellcolor[HTML]{FFFFFF}3.81E-09} & \multicolumn{1}{l|}{\cellcolor[HTML]{FFFFFF}6.05E-10} \\ \cline{2-7} 
\multicolumn{1}{l|}{}           & \multicolumn{1}{l|}{$2^{-4}$} & \multicolumn{1}{l|}{\cellcolor[HTML]{C0C0C0}1.76E-08} & \multicolumn{1}{l|}{\cellcolor[HTML]{C0C0C0}3.47E-09} & \multicolumn{1}{l|}{\cellcolor[HTML]{C0C0C0}9.74E-10} & \multicolumn{1}{l|}{\cellcolor[HTML]{C0C0C0}1.67E-10} & \multicolumn{1}{l|}{\cellcolor[HTML]{C0C0C0}2.45E-11} \\ \cline{2-7} 
\multicolumn{1}{l|}{$\Delta t$} & \multicolumn{1}{l|}{$2^{-5}$} & \multicolumn{1}{l|}{\cellcolor[HTML]{FFFFFF}9.59E-10} & \multicolumn{1}{l|}{\cellcolor[HTML]{FFFFFF}1.74E-10} & \multicolumn{1}{l|}{\cellcolor[HTML]{FFFFFF}4.71E-11} & \multicolumn{1}{l|}{\cellcolor[HTML]{FFFFFF}9.79E-12} & \multicolumn{1}{l|}{\cellcolor[HTML]{FFFFFF}2.80E-12} \\ \cline{2-7} 
\multicolumn{1}{l|}{}           & \multicolumn{1}{l|}{$2^{-6}$} & \multicolumn{1}{l|}{\cellcolor[HTML]{C0C0C0}7.55E-11} & \multicolumn{1}{l|}{\cellcolor[HTML]{C0C0C0}9.29E-12} & \multicolumn{1}{l|}{\cellcolor[HTML]{C0C0C0}3.27E-12} & \multicolumn{1}{l|}{\cellcolor[HTML]{C0C0C0}2.41E-12} & \multicolumn{1}{l|}{\cellcolor[HTML]{C0C0C0}1.96E-12} \\ \cline{2-7} 
\multicolumn{1}{l|}{}           & \multicolumn{1}{l|}{$2^{-7}$} & \multicolumn{1}{l|}{\cellcolor[HTML]{FFFFFF}1.47E-11} & \multicolumn{1}{l|}{\cellcolor[HTML]{FFFFFF}2.45E-12} & \multicolumn{1}{l|}{\cellcolor[HTML]{FFFFFF}2.86E-12} & \multicolumn{1}{l|}{\cellcolor[HTML]{FFFFFF}3.12E-12} & \multicolumn{1}{l|}{\cellcolor[HTML]{FFFFFF}2.01E-12} \\ \cline{2-7} 
\multicolumn{1}{l|}{}           & \multicolumn{1}{l|}{$2^{-8}$} & \multicolumn{1}{l|}{\cellcolor[HTML]{C0C0C0}6.00E-12} & \multicolumn{1}{l|}{\cellcolor[HTML]{C0C0C0}3.07E-12} & \multicolumn{1}{l|}{\cellcolor[HTML]{C0C0C0}3.12E-12} & \multicolumn{1}{l|}{\cellcolor[HTML]{C0C0C0}3.55E-12} & \multicolumn{1}{l|}{\cellcolor[HTML]{C0C0C0}2.21E-12} \\ \cline{2-7} 
\end{tabular}
\caption{The values of $\epsilon_{\alpha}$ for $\alpha=16/1024,...,64/1024$.}
\label{tbl:alpha:energy}
\end{table}

\begin{table}[!htb]
    \begin{subtable}{.5\linewidth}
      \centering
        \begin{tabular}{llll}
                                &                               & $\alpha$                                              &                                                       \\ \cline{3-4} 
                                & \multicolumn{1}{l|}{}         & \multicolumn{1}{l|}{8/1024}                           & \multicolumn{1}{l|}{12/1024}                          \\ \cline{2-4} 
\multicolumn{1}{l|}{}           & \multicolumn{1}{l|}{$2^{-2}$} & \multicolumn{1}{l|}{\cellcolor[HTML]{C0C0C0}3.82E-04} & \multicolumn{1}{l|}{\cellcolor[HTML]{C0C0C0}3.35E-05} \\ \cline{2-4} 
\multicolumn{1}{l|}{}           & \multicolumn{1}{l|}{$2^{-3}$} & \multicolumn{1}{l|}{6.91E-05}                         & \multicolumn{1}{l|}{1.60E-06}                         \\ \cline{2-4} 
\multicolumn{1}{l|}{}           & \multicolumn{1}{l|}{$2^{-4}$} & \multicolumn{1}{l|}{\cellcolor[HTML]{C0C0C0}3.05E-05} & \multicolumn{1}{l|}{\cellcolor[HTML]{C0C0C0}1.26E-07} \\ \cline{2-4} 
\multicolumn{1}{l|}{$\Delta t$} & \multicolumn{1}{l|}{$2^{-5}$} & \multicolumn{1}{l|}{1.57E-05}                         & \multicolumn{1}{l|}{3.14E-08}                         \\ \cline{2-4} 
\multicolumn{1}{l|}{}           & \multicolumn{1}{l|}{$2^{-6}$} & \multicolumn{1}{l|}{\cellcolor[HTML]{C0C0C0}8.05E-06} & \multicolumn{1}{l|}{\cellcolor[HTML]{C0C0C0}1.41E-08} \\ \cline{2-4} 
\multicolumn{1}{l|}{}           & \multicolumn{1}{l|}{$2^{-7}$} & \multicolumn{1}{l|}{4.08E-06}                         & \multicolumn{1}{l|}{6.99E-09}                         \\ \cline{2-4} 
\multicolumn{1}{l|}{}           & \multicolumn{1}{l|}{$2^{-8}$} & \multicolumn{1}{l|}{\cellcolor[HTML]{C0C0C0}2.05E-06} & \multicolumn{1}{l|}{\cellcolor[HTML]{C0C0C0}3.50E-09} \\ \cline{2-4} 
\end{tabular}
\caption{$N=128^3$}
    \end{subtable}%
    \begin{subtable}{.5\linewidth}
      \centering
        \begin{tabular}{llll}
                                &                               & $\alpha$                                              &                                                       \\ \cline{3-4} 
                                & \multicolumn{1}{l|}{}         & \multicolumn{1}{l|}{8/1024}                           & \multicolumn{1}{l|}{12/1024}                          \\ \cline{2-4} 
\multicolumn{1}{l|}{}           & \multicolumn{1}{l|}{$2^{-2}$} & \multicolumn{1}{l|}{\cellcolor[HTML]{C0C0C0}2.91E-04} & \multicolumn{1}{l|}{\cellcolor[HTML]{C0C0C0}3.33E-05} \\ \cline{2-4} 
\multicolumn{1}{l|}{}           & \multicolumn{1}{l|}{$2^{-3}$} & \multicolumn{1}{l|}{1.49E-05}                         & \multicolumn{1}{l|}{1.50E-06}                         \\ \cline{2-4} 
\multicolumn{1}{l|}{}           & \multicolumn{1}{l|}{$2^{-4}$} & \multicolumn{1}{l|}{\cellcolor[HTML]{C0C0C0}6.85E-07} & \multicolumn{1}{l|}{\cellcolor[HTML]{C0C0C0}7.21E-08} \\ \cline{2-4} 
\multicolumn{1}{l|}{$\Delta t$} & \multicolumn{1}{l|}{$2^{-5}$} & \multicolumn{1}{l|}{3.40E-08}                         & \multicolumn{1}{l|}{3.80E-09}                         \\ \cline{2-4} 
\multicolumn{1}{l|}{}           & \multicolumn{1}{l|}{$2^{-6}$} & \multicolumn{1}{l|}{\cellcolor[HTML]{C0C0C0}1.89E-09} & \multicolumn{1}{l|}{\cellcolor[HTML]{C0C0C0}2.14E-10} \\ \cline{2-4} 
\multicolumn{1}{l|}{}           & \multicolumn{1}{l|}{$2^{-7}$} & \multicolumn{1}{l|}{1.38E-10}                         & \multicolumn{1}{l|}{1.20E-11}                         \\ \cline{2-4} 
\multicolumn{1}{l|}{}           & \multicolumn{1}{l|}{$2^{-8}$} & \multicolumn{1}{l|}{\cellcolor[HTML]{C0C0C0}2.37E-11} & \multicolumn{1}{l|}{\cellcolor[HTML]{C0C0C0}2.80E-12} \\ \cline{2-4} 
\end{tabular}
\caption{$N=256^3, 512^3, 1024^3$}
    \end{subtable} 
    \caption{The values of $\epsilon_{\alpha}$ for $\alpha=8/1024,12/1024$.}
    \label{tbl:alpha:energy:2}
\end{table}

\section{Energy Spectra and Validation of Numerical Resolution}

The energy spectrum of the solution of (\ref{eq:Euler:Voigt}) is defined as,
\begin{equation}\label{eq:spectrum}
e(k,t) := \frac{1}{2}\sum_{j\leq|\bds j|<j+1}|\widehat{\u}_{\bds j}(t)|^2, \quad k=2\pi j, \quad \bds j\in\mathbb{Z}^3.
\end{equation}
The regularization term in (\ref{eq:Euler:Voigt}) acts as a filter which suppresses the growth of small-scale structures with length-scales smaller than the regularization parameter $\alpha$ \cite{molfetta_self-truncation_2015}. Due to this property, the solution to (\ref{eq:Euler:Voigt}) can remain well-resolved with relatively low resolution models. This is ideal for our analysis because it ensures that the model can fully capture the small scale structures of the fluid and so we can be certain that our objective functional (\ref{eq:objective:functional}) is accurately approximated.
\\~\\
{\color{red}
We know that for smooth solutions with Gevrey regularity, cf. (\ref{eq:Gevrey:Class}), the energy spectrum (\ref{eq:spectrum}) decays exponentially as the wavenumber $j=|\j|$ increases. Thus, whether or not a solution is well resolved can be assessed by examining the energy spectrum. Figures \ref{fig:spectrum:underresolved} and \ref{fig:spectrum:resolved} show the energy spectrum of the solution corresponding to $\be_{S}$ with $\alpha=8/1024$ and $\be_{TGV}$ with $\alpha=64/1024$, respectively. The spectrum is plotted over the time interval $T=5$. These figures demonstrate the exponential decay in the Fourier coefficients expected of a solution with Gevrey regularity. As the solutions evolve in time the rate of decay decreases as energy is transferred to small-scale structures. In Figure \ref{fig:spectrum:underresolved} a portion of the spectrum is cut-off by the 2/3 dealiasing filter. Due to the regularization term in (\ref{eq:Euler:Voigt}), the solution in Figure \ref{fig:spectrum:resolved} remains well resolved. }

\begin{figure}[!h]
\centering
\includegraphics[scale=0.7]{"figures/spectrum_underresolved.png"}
\caption{$e(k,t)$ vs. $k$ starting from $\be_S$ for $\alpha=8/1024$ and $N=256^3$ at $t=0$ (orange), $t=1$ (cyan), $t=2$ (brown), $t=3$ (blue), $t=4$ (green), $t=5$ (purple).}
\label{fig:spectrum:underresolved}
\end{figure}

\begin{figure}[!h]
\centering
\includegraphics[scale=0.7]{"figures/spectrum_resolved.png"}
\caption{$e(k,t)$ vs. $k$ starting from $\be_{TGV}$ for $\alpha=64/1024$ and $N=128^3$ at $t=0$ (orange), $t=1$ (cyan), $t=2$ (brown), $t=3$ (blue), $t=4$ (green), $t=5$ (purple).}
\label{fig:spectrum:resolved}
\end{figure}

\noindent {\color{red}In order to minimize the computational cost, the optimization searches were started from the initial guess at the lowest resolution, and ran until condition (\ref{eq:stop}) was satisfied. If the flow over $[0,T]$ corresponding to the optimal initial condition $\beat$ was under-resolved, then the initial condition was up-scaled to the next higher resolution $2N$ by padding with zeros in Fourier space and then converting back to physical space. The optimization search was restarted, now with the up-scaled $\beat$ as the initial guess $\beza$. Optimization was stopped when criterion (\ref{eq:stop}) was satisfied and the flow corresponding to $\tilde{\be}_{\alpha,T}$ remained well resolved, as verified by inspecting the energy spectrum (\ref{eq:spectrum}). Using this method we can ensure that the final optimal initial conditions and the corresponding flows were all well-resolved, and so we can be certain our evaluation of the objective functional is accurate for all of the cases discussed in Chapter \ref{chapter:5}.}

\section{Validation of the Gradient of the Objective Functional (\ref{eq:objective:functional})}

In the previous sections we examined the accuracy of the objective functional (\ref{eq:objective:functional}) and $\alpha$-energy (\ref{eq:alpha:Energy}) based on the regularization parameter $\alpha$, resolution $N$, and time step size $\Delta t$. In order to assess how accurately the gradient (\ref{eq:gradient}) of the objective functional is evaluated, we introduce the quantity \cite{ayala_extreme_2017,zhao_systematic_2023}. 
\begin{equation}\label{eq:kappa}
\kappa(\epsilon) = \frac{\epsilon^{-1}\left[ \pat(\bea+\epsilon\bepa)-\pat(\bea) \right]}{\la \nabla\pat(\bea),\bepa \ra_{G^\sigma}},
\end{equation}
where the numerator is the first-order finite difference approximation of the G\^{a}teaux differential (\ref{eq:Gateaux:differential}), and the denominator is the Reisz form (\ref{eq:Gateaux:differential:riesz}) of the differential.
\\~\\
A log-log plot of $|\kappa-1|$ vs $\epsilon$ can be broken up into three regions, which are dominated by roundoff errors, truncation errors, and discretization errors. The roundoff errors are caused by subtractive cancellation errors due to finite-precision arithmetics in the numerator of {\ref{eq:kappa}} and so are approximately f order $\O(\epsilon^{-1})$. Therefore, (\ref{eq:kappa}) will be dominated by round off errors for small $\epsilon$ values. The truncation errors are caused by the finite difference formula to approximate the G\^{a}teaux differential in the numerator and are of order $\O(\epsilon)$. Therefore, (\ref{eq:kappa}) will be dominated by the truncation errors for large values of $\epsilon$. Between these two regions (\ref{eq:kappa}) will be dominated by the discretization errors in the approximation of the PDEs, which depend on the spatial resolution $N$ and time step size $\Delta t$. Refining the discretization in both space and time ill reduce these errors. Therefore, if the calculation of the gradient is accurate, the log-log plot of $|\kappa-1|$ vs $\epsilon$ will be of the general form shown in Figure \ref{fig:kappa:theory}. 

\begin{figure}[h!]
\centering
\includegraphics[scale=0.9]{"figures/KappaTheory2.png"}
\caption{Schematic illustration of the log-log dependence of $\kappa(\epsilon)$ on $\epsilon$, cf. (\ref{eq:kappa}).}
\label{fig:kappa:theory}
\end{figure}

\noindent The quantity (\ref{eq:kappa}) was evaluated by fixing the spacial resolution at $N=512^3$, to keep the model well resolved, and varying the time step size $\Delta t \in [2^{-2},2^{-8}]$. The simulations were run for the time window of $T=5$ and $\epsilon\in [\text{1E-14}, 1]$. The values for $\alpha=8/1024$ and $\alpha=64/1024$ are shown in Figures \ref{fig:kappa:8} and \ref{fig:kappa:64}, respectively. The results from the $\kappa$ test reveal behaviour consistent with the expected behaviour illustrated in Figure \ref{fig:kappa:theory}. In particular, the values of $\kappa(\epsilon)$ obtained for intermediate values of $\epsilon$, the plateaux in the figures, approach unity as the time step is refined. Therefore, we conclude that the numerical techniques used to compute the gradient are accurate.

\begin{figure}[!h]
\centering
\includegraphics[scale=0.7]{"figures/kappa_8.png"}
\caption{Dependence of $\kappa(\epsilon)$ on $\epsilon$ for $\alpha=8/1024$. }
\label{fig:kappa:8}
\end{figure}

\begin{figure}[!h]
\centering
\includegraphics[scale=0.7]{"figures/kappa_64.png"}
\caption{Dependence of $\kappa(\epsilon)$ on $\epsilon$ for $\alpha=64/1024.$}
\label{fig:kappa:64}
\end{figure}

\section{Determination of the parameter defining the Gevrey space}\label{sec:sigma}

In Section \ref{sec:initial:conditions} we introduced the Gevrey space (\ref{eq:Gevrey:Class}) in which we search for optimal initial conditions. The relative 'size' of the space is determined by two parameters, $s\geq 3$ and $\sigma>0$. We use a fixed value of $s=3$. In order to determine an optimal value of $\sigma$, we solve the optimization problem for $\alpha=16/1024$, $T=5$, and different values of $\sigma = 10^{-1},...,10^{-5}$. The resulting values of the objective functional are compared in Figure \ref{fig:sigma}. 
\begin{figure}[!h]
\centering
\includegraphics[scale=0.7]{"figures/Sigma_2.png"}
\caption{$\pat(\be_{\alpha,T}^{(n)})$ vs iteration number $n$ for $\alpha=16/1024$, $T=5$, and different indicated values of $\sigma$.}
\label{fig:sigma}
\end{figure}

\noindent The results of Figure \ref{fig:sigma} show that the number of iterations it took to reach the maximum value does change depending on the value of $\sigma$, but there was not a significant variation in the value of the objective functional. This was also the result of the tests conducted for different values of $\alpha$. These results are consistent with those of \cite{zhao_systematic_2023} which found that the optimization results did not significantly change so long as $\sigma < 10^{-1}$. Therefore, we make the same selection as in \cite{zhao_systematic_2023} and choose $\sigma = 10^{-5}$.


\pagebreak